% s00-front.tex —— 前置件：译本说明
\chapter*{译本说明}
\addcontentsline{toc}{chapter}{译本说明}

\section*{一、底本}

本译本的底本是 Peter Scholze 于 2020 年秋季至 2021 年春季在 Berkeley 讲授的课程讲义的整理稿：

\begin{quote}
Peter Scholze（讲授），Tony Feng（记录），\emph{Geometrization of the Local Langlands
Correspondence}，2022 年 3 月 19 日生成的 PDF，共 148 页，27 讲附参考文献。
课程自 2020 年 10 月 30 日开始，至 2021 年 2 月 12 日结束；讲次标题后括注的月日即该讲日期。
\end{quote}

底本的印刷页码与 PDF 页码一致（偏移 0）。译本每一页的第一个段落之前都标有原书页码，
读者可据此与底本逐页对照。

\section*{二、体例}

\begin{enumerate}
  \item 术语以本仓库「数译规范」已冻结的译名为准（其中 \emph{perfectoid space} 译「完满空间」，
        \emph{diamond} 译「菱形」）；本库未收录者，参照香蕉空间（\texttt{bananaspace.org}）
        与通行中译取一译名，并在术语表中登记。
  \item 一个英文术语在全书中只对应一个中文译名；术语首次出现时以「中文（English）」形式给出原文。
  \item 人名、定理名、软件名、记号（如 $\Bun_G$、${\Bun_G}$、shtuka 等）不译。
        中文与拉丁字母之间留一个半角空格。
  \item 原文的条目编号（定义、定理、例、注、警告等）由 \LaTeX{} 共用计数器自动生成，
        与底本逐条对应；译文不手写编号。
  \item 公式一律依照底本页面图像重排，不依文本层转写的结果。底本公式偶有排版疏漏，
        凡影响文意者，以译注（脚注形式，标注等级与出处）说明，不改动正文。
\end{enumerate}

\section*{三、说明}

本篇为学术文献的中文译注本，供研读之用；原文著作权归原作者与记录者所有。
底本为作者公开在个人主页上的讲义 PDF，本译本不复制其页面图像，只重排文字与公式。
译文中若有错误，均属译者理解之误，与原作者无关。
